
\paragraph{Preguntas para investigar}
% Con las siguientes preguntas, trabajaremos en comprender definiciones.
\begin{enumerate}[label=P\arabic*]
	\item\label{item:perfectos}
		?`Qu\'e es un \emph{n\'umero perfecto}?
		Encontrar cuatro ejemplos.
	\item\label{item:mersenne}
		?`Qu\'e es un \emph{primo de Mersenne}?
		Encontrar cuatro ejemplos.
	\item\label{item:divisores}
		?`Qu\'e significa la expresi\'on
		\begin{math}
			\sum_{d\mid N,\,d>0}\,\frac 1 d
		\end{math}?
		Calcular el valor de esta sumatoria para algunos valores
		chicos de $N$.
	\seti
\end{enumerate}
%

\paragraph{Una pregunta para reflexionar}
Podemos nombrar varios sistemas matem\'aticos, cada uno con sus operaciones.
\begin{enumerate}[label=P\arabic*]
	\conti
	\item\label{item:sistemas}
		Compar\'a los siguientes sistemas matem\'aticos:
		$\Enteros$, $\Racionales$, $\Reales$, $2\Enteros$,
		$\Enterosmod[3]$, $\Enterosmod[6]$, $\Enterosmod[8]$,
		$\Enterosmod[11]$. Cada uno de ellos posee dos
		operaciones: suma y multiplicaci\'on ?`Cu\'ales de
		ellos se parecen teniendo en cuenta las propiedades
		de sus operaciones?
	\seti
\end{enumerate}
%

\paragraph{Una pregunta para explorar}
Podemos construir nuevos sistemas. Consideremos el conjunto de los
polinomios en $\indet$ con coeficientes en $\Enterosmod[3]$; denotamos
este conjunto por $\polinomios{\Enterosmod[3]}$.
\begin{enumerate}[label=P\arabic*]
	\conti
	\item\label{item:polinomios}
		Las expresiones $f(\indet)=\indet^2+2\indet+1$ y
		$g(\indet)=2\indet^2+\indet+1$ pertenecen a
		$\polinomios{\Enterosmod[3]}$. Calcul\'a
		$f(\indet)+g(\indet)$ y $f(\indet)\cdot g(\indet)$.
	\item\label{item:polinomios:factorizar}
		Factoriz\'a $\indet^4-1$ como producto de polinomios
		lineales en $\polinomios{\Enterosmod[5]}$.
		Factoriz\'a $\indet^6-1$ como producto de polinomios
		lineales en $\polinomios{\Enterosmod[7]}$.
	\seti
\end{enumerate}
%

\paragraph{Propiedades de los n\'umeros enteros}
Para responder algunas preguntas planteadas en estos problemas,
debemos recurrir a varias propiedades que comunmente asociamos
con los n\'umeros enteros. Hac\'e una lista de tales propiedades
y determin\'a si esa lista es suficiente para discutir cada pregunta
sobre aritm\'etica. La lista podr\'a ir cambiando a medida que
nos formulemos nuevas preguntas. Si la lista es adecuada,
podremos considerar que es una descripci\'on aceptable de los
n\'umeros enteros.

A continuaci\'on consideramos algunas propiedades de la divisibilidad.
Demostr\'a o refut\'a cada una
?`Es posible enmendar las afirmaciones falsas para que sean v\'alidas?
\begin{enumerate}[label=P\arabic*]
	\conti
	\item\label{item:divisibilidad:reflexiva}
		Si $a\in\Enteros$, entonces $a\mid a$.
	\item\label{item:divisibilidad:simetrica}
		Si $a,b\in\Enteros$ y $a\mid b$, entonces $b\mid a$.
	\item\label{item:divisibilidad:transitiva}
		Si $a,b,c\in\Enteros$ y $a\mid b$ y $b\mid c$,
		entonces $a\mid c$.
	\item\label{item:divisibilidad:multiplos}
		Si $a,b,c\in\Enteros$ y $a\mid b$, entonces $a\mid bc$.
	\item\label{item:divisibilidad:sumas}
		Si $a,b,c\in\Enteros$ y $a\mid b$ y $a\mid c$,
		entonces $a\mid (b+c)$.
	\item\label{item:divisibilidad:primos}
		Si $a,b,c\in\Enteros$ y $a\mid bc$, entonces
		$a\mid b$ o $a\mid c$.
	\seti
\end{enumerate}
%

\paragraph{Problemas num\'ericos}
Sigamos trabajando con los conceptos que tenemos.
\begin{enumerate}[label=P\arabic*]
	\conti
	\item\label{item:bases:cambiar}
		Usando divisi\'on en base siete, escrib\'{\i}
		$N=(34652)_7$ en las bases dos, cinco y once.
	\item\label{item:bases:escribir}
		Escrib\'{\i} cada uno de los siguientes n\'umeros
		en base tres: $3$, $9$, $27$, $243$, $1/3$, $1/9$
		y $1/27$. Escribilos en base dos ?`Qu\'e pod\'es
		decir?
	\item\label{item:modulares:varios}
		Hall\'a los siguientes elementos en $\Enterosmod[5]$:
		$-1$, $1/2$, $2/3$, $\sqrt{-1}$ ?`Cu\'antos de
		\'estos pod\'es encontrar en $\Enterosmod[6]$?
		?`Cu\'antos en $\Enterosmod[10]$, en $\Enterosmod[11]$
		y en $\Enterosmod[13]$?
	\item\label{item:modulares:cuadrados}
		Hac\'e una lista de los cuadrados perfectos en cada uno
		de los siguientes sistemas: $\Enterosmod[5]$,
		$\Enterosmod[17]$, $\Enterosmod[19]$,
		$\Enterosmod[21]$ ?`Cu\'antos cuadrados hay en cada
		caso? ?`Ten\'es alguna conjetura?
	\item\label{item:divisores:perfectos}
		Calcul\'a los valores de las siguientes sumatorias:
		\begin{math}
			\sum_{d\mid 6,\,d>0}\,\frac 1 d
		\end{math},
		\begin{math}
			\sum_{d\mid 28,\,d>0}\,\frac 1 d
		\end{math}
		?`Pod\'es hacer alguna conjetura?
		?`Si agregamos un ejemplo m\'as?
		Calcul\'a
		\begin{math}
			\sum_{d\mid 496,\,d>0}\,\frac 1 d
		\end{math}
		y compar\'a su valor con los valores de las sumatorias
		anteriores.
	\seti
\end{enumerate}
%

\paragraph{El arte de contar}
Un \emph{conjunto} est\'a conformado por \emph{elementos}.
Si nos quedamos s\'olo con algunos de los elementos de un conjunto,
tenemos un \emph{subconjunto}.
\begin{enumerate}[label=P\arabic*]
	\conti
	\item\label{item:subconjuntos}
		Si $S=\{a,b,c,d\}$, ?`cu\'antos subconjuntos de
		$S$ distintos hay?
		Si $T=\{c,d,e,f,g,h\}$, ?`cu\'antos hay en este caso?
	\seti
\end{enumerate}
%
Los conjuntos $S$ y $T$ tienen cuatro y seis elementos, respectivamente
\begin{enumerate}[label=P\arabic*]
	\conti
	\item\label{item:conjuntos:operaciones}
		?`Cu\'antos elementos hay en $S\cap T$?
		?`Cu\'antos elementos hay en $S\cup T$?
	\item\label{item:conjuntos:operaciones:bis}
		En general, si $S$ y $T$ son conjuntos cualesquiera y
		sabemos cu\'antos elementos tiene $S$ y cu\'antos
		tiene $T$, ?`podemos saber cu\'antos elementos tienen
		$S\cap T$ y $S\cup T$?
\end{enumerate}
%

